% Tables are pasted from the R console. Compile twice for the contents page.
\documentclass[12pt,letterpaper]{article}
\usepackage[T1]{fontenc}
\usepackage{lmodern}
\usepackage{graphicx,textcomp}
\usepackage{natbib}
\usepackage{setspace}
\usepackage{fullpage}
\usepackage{color}
\usepackage[reqno]{amsmath}
\usepackage{amsthm}
\usepackage{fancyvrb}
\usepackage{amssymb,enumerate}
\usepackage[all]{xy}
\usepackage{endnotes}
\usepackage{lscape}
\newtheorem{com}{Comment}
\usepackage{float}
\usepackage[hidelinks]{hyperref}
\newtheorem{lem} {Lemma}
\newtheorem{prop}{Proposition}
\newtheorem{thm}{Theorem}
\newtheorem{defn}{Definition}
\newtheorem{cor}{Corollary}
\newtheorem{obs}{Observation}
\usepackage[compact]{titlesec}
\usepackage{dcolumn}
\usepackage{tikz}
\usetikzlibrary{arrows}
\usepackage{multirow}
\usepackage{xcolor}
\newcolumntype{.}{D{.}{.}{-1}}
\newcolumntype{d}[1]{D{.}{.}{#1}}
\definecolor{light-gray}{gray}{0.65}
\usepackage{url}
\usepackage{listings}
\usepackage{color}

\definecolor{codegreen}{rgb}{0,0.6,0}
\definecolor{codegray}{rgb}{0.5,0.5,0.5}
\definecolor{codepurple}{rgb}{0.58,0,0.82}
\definecolor{backcolour}{rgb}{0.95,0.95,0.92}

\lstdefinestyle{mystyle}{
	backgroundcolor=\color{backcolour},   
	commentstyle=\color{codegreen},
	keywordstyle=\color{magenta},
	numberstyle=\tiny\color{codegray},
	stringstyle=\color{codepurple},
	basicstyle=\fontsize{10}{11}\selectfont,
	breakatwhitespace=false,         
	breaklines=true,                 
	captionpos=b,                    
	keepspaces=true,                 
	numbers=left,                    
	numbersep=5pt,                  
	showspaces=false,                
	showstringspaces=false,
	showtabs=false,                  
	tabsize=2
}
\lstset{style=mystyle,columns=fullflexible,upquote=true,aboveskip=2pt,belowskip=2pt}
\newcommand{\Sref}[1]{Section~\ref{#1}}
\newtheorem{hyp}{Hypothesis}

\setcounter{secnumdepth}{0}
\addtolength{\textheight}{16pt}
\addtolength{\footskip}{8pt}
\raggedbottom

% Set new command to make easier inserting Rcode.
\newcommand{\rcode}[2]{\lstinputlisting[language=R,firstline=#1,lastline=#2,firstnumber=1]{PS03_DBM.R}}

% Set new command to not have to repeat code for Answer: in bold and vspaces
\newcommand{\answer}{\par\vspace{.18cm}\noindent\textbf{Answer:} }

% Set new command to make easier inserting Graphics.
\newcommand{\fitplot}[1]{\par\smallskip\centerline{\includegraphics[width=0.82\linewidth]{#1}}\smallskip}

\renewcommand{\widehat}[1]{\mathord{\tikz[baseline=(outcome.base)]{
  \node[inner sep=0pt,outer sep=0pt] (outcome) {$#1$};
  \draw[line width=.48pt,line join=round]
    ([yshift=1.1pt]outcome.north west) -- ([yshift=4pt]outcome.north) --
    ([yshift=1.1pt]outcome.north east);}}}

\title{Problem Set 3}
\author{Barrera Munoz Diego}
\date{Due: November 12, 2026}

\begin{document}
\setlength{\abovedisplayskip}{7pt plus 2pt minus 1pt}
\setlength{\belowdisplayskip}{7pt plus 2pt minus 1pt}
\maketitle
\tableofcontents

\newpage

\phantomsection
\section{Question 1}
\noindent We are interested in knowing how the difference in campaign spending between incumbent and challenger affects the incumbent's vote share.

\begin{enumerate}[(1)]
\setlength{\itemsep}{4pt}\setlength{\parsep}{2pt}
\item Run a regression where the outcome variable is \texttt{voteshare} and the explanatory variable is \texttt{difflog}.

\answer First, we import the data and check the three variables for missing data.

\rcode{12}{19}

% Pasted stargazer output
\par\nobreak
{\footnotesize\setlength{\tabcolsep}{4pt}\renewcommand{\arraystretch}{0.93}%
\noindent\makebox[\linewidth][c]{%
\begin{tabular}{@{\extracolsep{5pt}} ccccccccc} 
\\[-1.8ex]\hline 
\hline \\[-1.8ex] 
 & Elections & Missing & Min. & 1st Qu. & Median & Mean & 3rd Qu. & Max. \\ 
\hline \\[-1.8ex] 
voteshare & $3,193$ & $0$ & $0.3476$ & $0.5846$ & $0.6569$ & $0.6552$ & $0.7233$ & $0.9930$ \\ 
difflog & $3,193$ & $0$ & $$-$3.0601$ & $0.6734$ & $1.6106$ & $1.8278$ & $2.9857$ & $5.8558$ \\ 
presvote & $3,193$ & $0$ & $0.1931$ & $0.4695$ & $0.5478$ & $0.5512$ & $0.6240$ & $0.9606$ \\ 
\hline \\[-1.8ex] 
\end{tabular}
}\par}
\smallskip


All 3,193 elections have complete data (no NAs).

We regress incumbent \texttt{voteshare} $\sim$ \texttt{difflog} (log incumbent spending minus log challenger spending). The regression tables give standard errors in parentheses.

\rcode{22}{24}

% Pasted stargazer output
\par\nobreak
{\footnotesize\setlength{\tabcolsep}{4pt}\renewcommand{\arraystretch}{0.93}%
\noindent\makebox[\linewidth][c]{%
\begin{tabular}{@{\extracolsep{5pt}}lc} 
\\[-1.8ex]\hline 
\hline \\[-1.8ex] 
 & \multicolumn{1}{c}{\textit{Dependent variable:}} \\ 
\cline{2-2} 
\\[-1.8ex] & voteshare \\ 
\hline \\[-1.8ex] 
 difflog & 0.041666$^{***}$ (0.000968) \\ 
  & t = 43.044060 \\ 
  Constant & 0.579031$^{***}$ (0.002251) \\ 
  & t = 257.188300 \\ 
 \hline \\[-1.8ex] 
Observations & 3,193 \\ 
R$^{2}$ & 0.367341 \\ 
Adjusted R$^{2}$ & 0.367143 \\ 
\hline 
\hline \\[-1.8ex] 
\textit{Note:}  & \multicolumn{1}{r}{$^{*}$p$<$0.1; $^{**}$p$<$0.05; $^{***}$p$<$0.01} \\ 
\end{tabular}
}\par}
\smallskip


The spending gap explains $36.73\%$ of the variation in incumbent vote share. The slope is positive. Our null hypothesis is that the population spending-gap slope is zero. Our alternative
is that it differs from zero. The slope is positive ($t = 43.04$, two-sided $p < 0.001$), so
we reject the null at $5\%$. 

\newpage

\item Make a scatterplot of the two variables and add the regression line.

\answer First, we calculate the fitted mean and its pointwise $95\%$ confidence limits across
the observed spending gaps (to use later in the scatterplot). 

\rcode{29}{34}

\rcode{37}{50}

\fitplot{PS03_Q1_vote_spending.pdf}

Incumbents tend to receive a larger vote share as their spending advantage over the challenger grows. The spending gap explains $36.73\%$ of the variation in incumbent vote share. Individual election results vary much more widely around the fitted line than the dashed 95\% confidence limits for the mean.

\newpage

\item Save the residuals of the model in a separate object.

\answer We save the residuals and check the first election. A residual is the observed vote share minus the fitted vote share.

\rcode{55}{59}

% Pasted stargazer output
\par\nobreak
{\footnotesize\setlength{\tabcolsep}{4pt}\renewcommand{\arraystretch}{0.93}%
\noindent\makebox[\linewidth][c]{%
\begin{tabular}{@{\extracolsep{5pt}} ccc} 
\\[-1.8ex]\hline 
\hline \\[-1.8ex] 
observed & fitted & residual \\ 
\hline \\[-1.8ex] 
$0.602361$ & $0.602784$ & $$-$0.000423$ \\ 
\hline \\[-1.8ex] 
\end{tabular}
}\par}
\smallskip


The first incumbent received $60.24\%$ of the vote, about $0.042$ percentage points below the prediction. This explains the negative residual.

\item Write the prediction equation.

\answer We take the intercept and slope from the fitted model:

\rcode{64}{65}

% Pasted stargazer output
\par\nobreak
{\footnotesize\setlength{\tabcolsep}{4pt}\renewcommand{\arraystretch}{0.93}%
\noindent\makebox[\linewidth][c]{%
\begin{tabular}{@{\extracolsep{5pt}} cc} 
\\[-1.8ex]\hline 
\hline \\[-1.8ex] 
(Intercept) & difflog \\ 
\hline \\[-1.8ex] 
$0.579031$ & $0.041666$ \\ 
\hline \\[-1.8ex] 
\end{tabular}
}\par}
\smallskip


For election $i$, the prediction equation is:

$$\widehat{\texttt{voteshare}}_i=0.579031+0.041666\,\texttt{difflog}_i$$

On average, a one-unit change in the log spending gap in the incumbent's favour predicts a gain of $4.17$ percentage points in incumbent vote share, since the fitted slope is $\hat{\beta}_1=0.041666$. At equal spending ($\texttt{difflog}=0$), the intercept $\hat{\beta}_0=0.579031$ gives a predicted incumbent vote share of $57.90\%$.
\end{enumerate}

\newpage

\phantomsection
\section{Question 2}
\noindent We are interested in knowing how the difference between incumbent and challenger's spending and the vote share of the presidential candidate of the incumbent's party are related.

\begin{enumerate}[(1)]
\setlength{\itemsep}{4pt}\setlength{\parsep}{2pt}
\item Run a regression where the outcome variable is \texttt{presvote} and the explanatory variable is \texttt{difflog}.

\answer We will use \texttt{difflog} to predict the vote share of the presidential
candidate from the incumbent’s party (\texttt{presvote}).

\rcode{72}{74}

% Pasted stargazer output
\par\nobreak
{\footnotesize\setlength{\tabcolsep}{4pt}\renewcommand{\arraystretch}{0.93}%
\noindent\makebox[\linewidth][c]{%
\begin{tabular}{@{\extracolsep{5pt}}lc} 
\\[-1.8ex]\hline 
\hline \\[-1.8ex] 
 & \multicolumn{1}{c}{\textit{Dependent variable:}} \\ 
\cline{2-2} 
\\[-1.8ex] & presvote \\ 
\hline \\[-1.8ex] 
 difflog & 0.023837$^{***}$ (0.001359) \\ 
  & t = 17.541820 \\ 
  Constant & 0.507583$^{***}$ (0.003161) \\ 
  & t = 160.600800 \\ 
 \hline \\[-1.8ex] 
Observations & 3,193 \\ 
R$^{2}$ & 0.087951 \\ 
Adjusted R$^{2}$ & 0.087665 \\ 
\hline 
\hline \\[-1.8ex] 
\textit{Note:}  & \multicolumn{1}{r}{$^{*}$p$<$0.1; $^{**}$p$<$0.05; $^{***}$p$<$0.01} \\ 
\end{tabular}
}\par}
\smallskip


Our null hypothesis is that the population spending-gap slope is zero. Our two-sided
alternative is that it differs from zero. We reject the null at $5\%$ ($t = 17.54$, $p < 0.001$).
The association is positive. The gap explains $8.80\%$ of the variation in presidential
vote share.

\item Make a scatterplot of the two variables and add the regression line.

\answer We put \texttt{difflog} on the horizontal axis and \texttt{presvote} on the vertical axis. We calculate the 95\% confidence limits from the model in part~(1).

\rcode{79}{81}

\rcode{84}{97}

\fitplot{PS03_Q2_pres_spending.pdf}

\enlargethispage{\baselineskip}

The line slopes upward, but the points are more spread out than in Question~1. The spending gap tells us less about the presidential vote.

\item Save the residuals of the model in a separate object.

\answer We save the presidential residuals and check the first election:

\rcode{102}{106}

% Pasted stargazer output
\par\nobreak
{\footnotesize\setlength{\tabcolsep}{4pt}\renewcommand{\arraystretch}{0.93}%
\noindent\makebox[\linewidth][c]{%
\begin{tabular}{@{\extracolsep{5pt}} ccc} 
\\[-1.8ex]\hline 
\hline \\[-1.8ex] 
observed & fitted & residual \\ 
\hline \\[-1.8ex] 
$0.526778$ & $0.521173$ & $0.005606$ \\ 
\hline \\[-1.8ex] 
\end{tabular}
}\par}
\smallskip


The first residual is positive: the presidential candidate received $0.56$ percentage points more than predicted.

\item Write the prediction equation.

\answer We can read the two coefficients from the model:

\rcode{111}{112}

% Pasted stargazer output
\par\nobreak
{\footnotesize\setlength{\tabcolsep}{4pt}\renewcommand{\arraystretch}{0.93}%
\noindent\makebox[\linewidth][c]{%
\begin{tabular}{@{\extracolsep{5pt}} cc} 
\\[-1.8ex]\hline 
\hline \\[-1.8ex] 
(Intercept) & difflog \\ 
\hline \\[-1.8ex] 
$0.507583$ & $0.023837$ \\ 
\hline \\[-1.8ex] 
\end{tabular}
}\par}
\smallskip


The prediction equation is

$$\widehat{\texttt{presvote}}_i=0.507583+0.023837\,\texttt{difflog}_i$$

On average, a one-unit change in \texttt{difflog} from 0 to 1 predicts a $2.38$-percentage-point rise in the party's presidential vote share, since the slope is $\hat{\beta}_1=0.023837$. At equal spending, the intercept $\hat{\beta}_0=0.507583$ gives a predicted presidential vote share of $50.76\%$.
\end{enumerate}

\newpage

\phantomsection
\section{Question 3}
\noindent We are interested in knowing how the vote share of the presidential candidate of the incumbent's party is associated with the incumbent's electoral success.

\begin{enumerate}[(1)]
\setlength{\itemsep}{4pt}\setlength{\parsep}{2pt}
\item Run a regression where the outcome variable is \texttt{voteshare} and the explanatory variable is \texttt{presvote}.

\answer We regress \texttt{voteshare} $\sim$ \texttt{presvote}:

\rcode{119}{121}

% Pasted stargazer output
\par\nobreak
{\footnotesize\setlength{\tabcolsep}{4pt}\renewcommand{\arraystretch}{0.93}%
\noindent\makebox[\linewidth][c]{%
\begin{tabular}{@{\extracolsep{5pt}}lc} 
\\[-1.8ex]\hline 
\hline \\[-1.8ex] 
 & \multicolumn{1}{c}{\textit{Dependent variable:}} \\ 
\cline{2-2} 
\\[-1.8ex] & voteshare \\ 
\hline \\[-1.8ex] 
 presvote & 0.388018$^{***}$ (0.013493) \\ 
  & t = 28.756740 \\ 
  Constant & 0.441330$^{***}$ (0.007599) \\ 
  & t = 58.080330 \\ 
 \hline \\[-1.8ex] 
Observations & 3,193 \\ 
R$^{2}$ & 0.205814 \\ 
Adjusted R$^{2}$ & 0.205565 \\ 
\hline 
\hline \\[-1.8ex] 
\textit{Note:}  & \multicolumn{1}{r}{$^{*}$p$<$0.1; $^{**}$p$<$0.05; $^{***}$p$<$0.01} \\ 
\end{tabular}
}\par}
\smallskip


Our null hypothesis is that the population slope on presidential vote share is zero.
Our two-sided alternative is that it differs from zero. We reject the null at the $5\%$
level ($t = 28.76$, $p < 0.001$). Greater support for the party’s presidential candidate
predicts greater support for its incumbent. Presidential vote share explains $20.58\%$ of
the variation in incumbent vote share.

\item Make a scatterplot of the two variables and add the regression line.

\enlargethispage{\baselineskip}

\answer We plot one point per election and show both vote shares as percentages. The presidential vote share nearest $50\%$ in the grid for the 95\% confidence limits is $50.17\%$.

\rcode{126}{130}

\rcode{133}{147}

\fitplot{PS03_Q3_vote_pres.pdf}

Incumbents tend to receive a larger share of the vote where their party's presidential candidate does better. Even near a presidential vote share of $50\%$, we see incumbents who lost their race as well as incumbents who won comfortably.

\item Write the prediction equation.

\answer We use the two coefficients from the fitted line:

\rcode{152}{153}

% Pasted stargazer output
\par\nobreak
{\footnotesize\setlength{\tabcolsep}{4pt}\renewcommand{\arraystretch}{0.93}%
\noindent\makebox[\linewidth][c]{%
\begin{tabular}{@{\extracolsep{5pt}} cc} 
\\[-1.8ex]\hline 
\hline \\[-1.8ex] 
(Intercept) & presvote \\ 
\hline \\[-1.8ex] 
$0.441330$ & $0.388018$ \\ 
\hline \\[-1.8ex] 
\end{tabular}
}\par}
\smallskip


The prediction equation is

$$\widehat{\texttt{voteshare}}_i=0.441330+0.388018\,\texttt{presvote}_i$$

On average, a one-unit change in presidential vote share (100 percentage points) predicts a $38.80$-percentage-point change in incumbent vote share in the same direction, since the slope is $\hat{\beta}_1=0.388018$. A rise from $50\%$ to $60\%$ therefore predicts a gain of $0.388018(0.60-0.50)\approx0.03880$, or $3.88$ percentage points, for the incumbent. The intercept $\hat{\beta}_0=0.441330$ predicts $44.13\%$ incumbent vote share at zero presidential support. That baseline is an extrapolation because the smallest presidential vote share in the sample is $19.31\%$.

\end{enumerate}

\newpage

\phantomsection
\section{Question 4}
\noindent The residuals from part (a) tell us how much of the variation in \texttt{voteshare} is $not$ explained by the difference in spending between incumbent and challenger. The residuals in part (b) tell us how much of the variation in \texttt{presvote} is $not$ explained by the difference in spending between incumbent and challenger in the district.

\begin{enumerate}[(1)]
\setlength{\itemsep}{4pt}\setlength{\parsep}{2pt}
\item Run a regression where the outcome variable is the residuals from Question 1 and the explanatory variable is the residuals from Question 2.

\answer We can use the incumbent vote-share residuals saved in Question~1(3) as the outcome and the presidential vote-share residuals saved in Question~2(3) as the explanatory variable.

\rcode{160}{163}

% Pasted stargazer output
\par\nobreak
{\footnotesize\setlength{\tabcolsep}{4pt}\renewcommand{\arraystretch}{0.93}%
\noindent\makebox[\linewidth][c]{%
\begin{tabular}{@{\extracolsep{5pt}}lc} 
\\[-1.8ex]\hline 
\hline \\[-1.8ex] 
 & \multicolumn{1}{c}{\textit{Dependent variable:}} \\ 
\cline{2-2} 
\\[-1.8ex] & vote\_resid \\ 
\hline \\[-1.8ex] 
 pres\_resid & 0.256877$^{***}$ (0.011762) \\ 
  & t = 21.839750 \\ 
  Constant & $-$0.000000 (0.001299) \\ 
  & t = $-$0.000000 \\ 
 \hline \\[-1.8ex] 
Observations & 3,193 \\ 
R$^{2}$ & 0.130038 \\ 
Adjusted R$^{2}$ & 0.129765 \\ 
\hline 
\hline \\[-1.8ex] 
\textit{Note:}  & \multicolumn{1}{r}{$^{*}$p$<$0.1; $^{**}$p$<$0.05; $^{***}$p$<$0.01} \\ 
\end{tabular}
}\par}
\smallskip


Our null hypothesis is that the population slope relating the two spending-adjusted
shares is zero. Our two-sided alternative is that it differs from zero. We reject the null
at the $5\%$ level ($t = 21.84$, $p < 0.001$). The presidential residual explains $13.00\%$ of
the variation left in incumbent vote share after adjusting for spending.

\item Make a scatterplot of the two residuals and add the regression line.

\answer We put the presidential residual on the horizontal axis and the incumbent residual on the vertical axis.

\rcode{168}{172}

\rcode{175}{190}

\fitplot{PS03_Q4_added_variable.pdf}

Each point represents one election. The fitted line slopes upward. After adjusting both vote shares for spending, incumbents tend to do better than predicted where their party's presidential candidate also does better than predicted. The first election lies in the lower-right quadrant, with presidential vote share $0.56$ percentage points above its prediction and incumbent vote share $0.042$ percentage points below its prediction.

\item Write the prediction equation.

\answer We take the coefficients from the residual regression:

\rcode{195}{196}

% Pasted stargazer output
\par\nobreak
{\footnotesize\setlength{\tabcolsep}{4pt}\renewcommand{\arraystretch}{0.93}%
\noindent\makebox[\linewidth][c]{%
\begin{tabular}{@{\extracolsep{5pt}} cc} 
\\[-1.8ex]\hline 
\hline \\[-1.8ex] 
(Intercept) & pres\_resid \\ 
\hline \\[-1.8ex] 
$0$ & $0.256877$ \\ 
\hline \\[-1.8ex] 
\end{tabular}
}\par}
\smallskip


$$\widehat{\texttt{vote\_resid}}_i=0.256877\,\texttt{pres\_resid}_i$$

On average, a one-unit change in the presidential vote-share residual predicts a $0.256877$-unit change in the incumbent vote-share residual in the same direction, as measured by the slope $\hat{\beta}_1$. In percentage points, a presidential residual one point higher predicts an incumbent residual about $0.257$ points higher. Both sets of residuals average zero because their original regressions include intercepts. The fitted intercept $\hat{\beta}_0$ therefore rounds to zero, and the line passes through the origin.

\end{enumerate}

\newpage

\phantomsection
\section{Question 5}
\noindent What if the incumbent's vote share is affected by both the president's popularity and the difference in spending between incumbent and challenger?

\begin{enumerate}[(1)]
\setlength{\itemsep}{4pt}\setlength{\parsep}{2pt}
\item Run a regression where the outcome variable is the incumbent's \texttt{voteshare} and the explanatory variables are \texttt{difflog} and \texttt{presvote}.

\answer We will include the spending gap and presidential vote share in the same regression. Each slope measures the association with incumbent vote share while we hold the other predictor fixed.

\rcode{203}{205}

% Pasted stargazer output
\par\nobreak
{\footnotesize\setlength{\tabcolsep}{4pt}\renewcommand{\arraystretch}{0.93}%
\noindent\makebox[\linewidth][c]{%
\begin{tabular}{@{\extracolsep{5pt}}lc} 
\\[-1.8ex]\hline 
\hline \\[-1.8ex] 
 & \multicolumn{1}{c}{\textit{Dependent variable:}} \\ 
\cline{2-2} 
\\[-1.8ex] & voteshare \\ 
\hline \\[-1.8ex] 
 difflog & 0.035543$^{***}$ (0.000946) \\ 
  & t = 37.590140 \\ 
  presvote & 0.256877$^{***}$ (0.011764) \\ 
  & t = 21.836330 \\ 
  Constant & 0.448644$^{***}$ (0.006330) \\ 
  & t = 70.879480 \\ 
 \hline \\[-1.8ex] 
Observations & 3,193 \\ 
R$^{2}$ & 0.449610 \\ 
Adjusted R$^{2}$ & 0.449265 \\ 
\hline 
\hline \\[-1.8ex] 
\textit{Note:}  & \multicolumn{1}{r}{$^{*}$p$<$0.1; $^{**}$p$<$0.05; $^{***}$p$<$0.01} \\ 
\end{tabular}
}\par}
\smallskip


For each predictor, our null hypothesis is that its population slope is zero with the other predictor held fixed. The two-sided alternative is that its slope differs from zero. We reject both null hypotheses at the 5\% level $(p<0.001$ for each slope). Together the predictors explain $44.96\%$ of the variation in incumbent vote share.

\newpage

\item Write the prediction equation.

\answer We use the fitted intercept and both partial slopes:

\rcode{210}{211}

% Pasted stargazer output
\par\nobreak
{\footnotesize\setlength{\tabcolsep}{4pt}\renewcommand{\arraystretch}{0.93}%
\noindent\makebox[\linewidth][c]{%
\begin{tabular}{@{\extracolsep{5pt}} ccc} 
\\[-1.8ex]\hline 
\hline \\[-1.8ex] 
(Intercept) & difflog & presvote \\ 
\hline \\[-1.8ex] 
$0.448644$ & $0.035543$ & $0.256877$ \\ 
\hline \\[-1.8ex] 
\end{tabular}
}\par}
\smallskip


Therefore, the prediction equation is

$$\widehat{\texttt{voteshare}}_i=0.448644+0.035543\,\texttt{difflog}_i+0.256877\,\texttt{presvote}_i$$

At equal spending and zero presidential vote share, the intercept $\hat{\beta}_0=0.448644$ gives a predicted incumbent vote share of $44.86\%$. On average, a one-unit change in the log spending gap in the incumbent's favour predicts a gain of $3.55$ percentage points in incumbent vote share ($\hat{\beta}_1=0.035543$), holding presidential vote share constant. Holding the log spending gap constant, the slope $\hat{\beta}_2=0.256877$ gives the predicted gain in incumbent vote share for a one-unit rise in presidential vote share. A rise from $50\%$ to $60\%$ therefore predicts a $2.57$-percentage-point gain for the incumbent, compared with $3.88$ points in Question~3.

\item What is it in this output that is identical to the output in Question 4? Why do you think this is the case?

\answer The coefficient on \texttt{presvote} in Question~5 equals the coefficient on \texttt{pres\_resid} in Question~4 to numerical precision. The left column shows the residual regression from Question~4 and the right column shows the full regression from Question~5.

\rcode{216}{219}

% Pasted stargazer output
\par\nobreak
{\footnotesize\setlength{\tabcolsep}{4pt}\renewcommand{\arraystretch}{0.93}%
\noindent\makebox[\linewidth][c]{%
\begin{tabular}{@{\extracolsep{5pt}}lcc} 
\\[-1.8ex]\hline 
\hline \\[-1.8ex] 
 & \multicolumn{2}{c}{\textit{Dependent variable:}} \\ 
\cline{2-3} 
\\[-1.8ex] & vote\_resid & voteshare \\ 
\\[-1.8ex] & (1) & (2)\\ 
\hline \\[-1.8ex] 
 pres\_resid & 0.2568770127$^{***}$ (0.0117619024) &  \\ 
  difflog &  & 0.0355430864$^{***}$ (0.0009455428) \\ 
  presvote &  & 0.2568770127$^{***}$ (0.0117637458) \\ 
  Constant & $-$0.0000000000 (0.0012986050) & 0.4486442218$^{***}$ (0.0063296774) \\ 
 \hline \\[-1.8ex] 
Observations & 3,193 & 3,193 \\ 
R$^{2}$ & 0.1300376195 & 0.4496104332 \\ 
Adjusted R$^{2}$ & 0.1297649895 & 0.4492653614 \\ 
Residual Std. Error & 0.0733797985 (df = 3191) & 0.0733912991 (df = 3190) \\ 
\hline 
\hline \\[-1.8ex] 
\textit{Note:}  & \multicolumn{2}{r}{$^{*}$p$<$0.1; $^{**}$p$<$0.05; $^{***}$p$<$0.01} \\ 
\end{tabular}
}\par}
\smallskip

\newpage

We compare the unrounded presidential slopes first, then the residuals for all 3,193 elections.

\rcode{222}{225}

\begin{verbatim}
[1] TRUE
[1] TRUE
\end{verbatim}


Both slopes equal $0.2568770127$, and the residuals match to numerical precision. For a one-percentage-point increase in presidential vote share at the same spending gap, the model predicts about $0.257$ percentage points more incumbent vote share.

For the explanation of why this happens, we can see the slides~12--14 of the Week~9 tutorial:
 In Question~1, we subtract the incumbent vote share predicted by spending from the observed incumbent share. In Question~2, we do the same for the presidential share, using the same spending gap and the same elections. What remains in each share is uncorrelated with the spending gap. Regressing the first residual on the second therefore uses exactly the variation that estimates the presidential slope when we hold spending fixed in Question~5. This is the Frisch--Waugh--Lovell result.

The standard errors both round to $0.01176$. For the error variance, R divides the same residual sum of squares by $3{,}191$ in Question~4 and by $3{,}190$ in Question~5, where we estimate one more coefficient. This gives the slightly larger standard error in Question~5. The $R^2$ values use different totals as well. Question~4 accounts for $13.00\%$ of the incumbent variation remaining after spending, and Question~5 accounts for $44.96\%$ of the original variation.
\end{enumerate}
\end{document}
